\chapter{TỔNG QUAN VÀ PHÂN LOẠI CÁC NGHIÊN CỨU}
\label{chap:taxonomy}

\section{Khung phân loại bốn nhóm}

Dựa trên câu hỏi nghiên cứu chính mà mỗi bài giải quyết, nhóm phân loại tám nghiên cứu thành bốn hướng (Hình~\ref{fig:taxonomy}):

\begin{enumerate}
  \item \textbf{Safety \& Robustness} (\textcolor{gSafe}{\textbf{S}}): hệ thống có chống chịu được tấn công hay thông tin sai lệch không? Gồm MIRAGE (\#1) -- bảo vệ bí mật topology, và TIDE (\#4) -- đánh giá ảnh hưởng của thông tin sai đến quá trình khôi phục sự thật tập thể.
  \item \textbf{Orchestration \& Control} (\textcolor{gOrch}{\textbf{O}}): quyết định ai chạy, chạy bằng mô hình nào, và agent nào có rủi ro. Gồm ProgRouter (\#3) -- model routing theo tiến độ dưới ràng buộc chi phí, và HalluProp (\#5) -- suy luận rủi ro thất bại trước khi hệ thống chạy.
  \item \textbf{Communication \& Interaction} (\textcolor{gComm}{\textbf{C}}): các agent nên trao đổi thông tin gì, với ai và ở mức độ nào. Gồm DAIN (\#7) -- mạng tương tác động giữa các agent chuyên biệt cho dữ liệu đa phương thức, và AGENT-RADAR (\#8) -- attention steering theo mức độ liên quan của ngữ cảnh.
  \item \textbf{Domain Applications} (\textcolor{gApp}{\textbf{A}}): áp dụng multi-agent collaboration vào bài toán chuyên ngành. Gồm MAC-RRG (\#2) -- radiology report generation ngực, và UrbanAgent (\#6) -- dự đoán chỉ số kinh tế -- xã hội vùng đô thị.
\end{enumerate}

\begin{figure}[!htbp]
\centering
\resizebox{\linewidth}{!}{%
\begin{tikzpicture}[
  font=\small,
  root/.style={draw,rounded corners,fill=black!8,minimum width=5.2cm,minimum height=0.9cm,align=center,font=\bfseries\small},
  grp/.style={draw,rounded corners,minimum width=3.4cm,minimum height=0.9cm,align=center,font=\bfseries\footnotesize,text=white},
  leaf/.style={draw,rounded corners=1pt,minimum width=3.4cm,align=center,font=\footnotesize,fill=white},
  arr/.style={-{Stealth[length=2mm]},thick,black!60}
]
\node[root] (r) at (0,0) {Multi-Agent Systems \&\\Collaborative Reasoning};
\node[grp,fill=gSafe] (s) at (-6.0,-1.9) {Safety \&\\Robustness};
\node[grp,fill=gOrch] (o) at (-2.0,-1.9) {Orchestration\\\& Control};
\node[grp,fill=gComm] (c) at (2.0,-1.9)  {Communication\\\& Interaction};
\node[grp,fill=gApp]  (a) at (6.0,-1.9)  {Domain\\Applications};
\foreach \g in {s,o,c,a} \draw[arr] (r.south) -- (\g.north);

\node[leaf,below=0.45cm of s] (s1) {\#1 MIRAGE\\\scriptsize ẩn topology};
\node[leaf,below=0.2cm of s1] (s2) {\#4 TIDE\\\scriptsize thông tin sai \& sự thật phân tán};
\node[leaf,below=0.45cm of o] (o1) {\#3 ProgRouter\\\scriptsize routing theo tiến độ};
\node[leaf,below=0.2cm of o1] (o2) {\#5 HalluProp\\\scriptsize rủi ro trước khi chạy};
\node[leaf,below=0.45cm of c] (c1) {\#7 DAIN\\\scriptsize mạng tương tác động};
\node[leaf,below=0.2cm of c1] (c2) {\#8 AGENT-RADAR\\\scriptsize attention steering};
\node[leaf,below=0.45cm of a] (a1) {\#2 MAC-RRG\\\scriptsize báo cáo X-quang};
\node[leaf,below=0.2cm of a1] (a2) {\#6 UrbanAgent\\\scriptsize urban region profiling};
\foreach \g/\l in {s/s1,o/o1,c/c1,a/a1} \draw[arr] (\g) -- (\l);
\end{tikzpicture}}
\caption{Khung phân loại bốn nhóm cho tám nghiên cứu được khảo sát.}
\label{fig:taxonomy}
\end{figure}

\section{Ma trận multi-label}

Phân loại một nhãn chỉ phản ánh mục tiêu \emph{chính}. Trên thực tế, hầu hết các bài đều chạm tới nhiều nhóm: HalluProp được xếp vào điều phối vì đầu ra của nó dùng để cắt tỉa cạnh hoặc sửa vai trò trước khi chạy, nhưng bản chất là phương pháp an toàn (bài được nộp vào chuyên mục cs.CR); DAIN vừa là phương pháp giao tiếp vừa là phương pháp tiết kiệm tính toán qua sparse activation. Bảng~\ref{tab:crosslabel} thể hiện quan hệ đa nhãn này.

\begin{table}[!htbp]
\centering
\small
\caption{Ma trận multi-label: \textbf{\large$\bullet$} = hướng chính, $\circ$ = hướng phụ.}
\label{tab:crosslabel}
\begin{tabular}{l C{1.8cm} C{1.8cm} C{1.8cm} C{2.25cm}}
\toprule
\textbf{Nghiên cứu} & \textcolor{gSafe}{\textbf{Safety}} & \textcolor{gOrch}{\textbf{Orchestr.}} & \textcolor{gComm}{\textbf{Comm.}} & \textcolor{gApp}{\textbf{Application}} \\
\midrule
\#1 MIRAGE      & \textbf{\large$\bullet$} &          & $\circ$  &          \\
\#4 TIDE        & \textbf{\large$\bullet$} &          & $\circ$  &          \\
\#3 ProgRouter  &          & \textbf{\large$\bullet$} &          &          \\
\#5 HalluProp   & $\circ$  & \textbf{\large$\bullet$} & $\circ$  &          \\
\#7 DAIN        &          & $\circ$  & \textbf{\large$\bullet$} & $\circ$ (y tế) \\
\#8 AGENT-RADAR &          & $\circ$  & \textbf{\large$\bullet$} &          \\
\#2 MAC-RRG     &          & $\circ$  &          & \textbf{\large$\bullet$} \\
\#6 UrbanAgent  &          &          & $\circ$  & \textbf{\large$\bullet$} \\
\bottomrule
\end{tabular}
\end{table}

\section{Vị trí can thiệp trong lifecycle của MAS}

Một cách nhìn bổ sung là xét \emph{thời điểm} mỗi phương pháp tác động lên vòng đời của một \MAS{} (Hình~\ref{fig:lifecycle}):
\begin{itemize}
  \item \textbf{Trước khi chạy (pre-hoc):} HalluProp chỉ dùng role prompt, câu hỏi và đồ thị để ước lượng rủi ro, chưa có agent nào sinh ra câu trả lời.
  \item \textbf{Trong khi chạy -- điều phối:} ProgRouter chọn mô hình cho từng bước; DAIN chọn agent nào được kích hoạt cho từng mẫu.
  \item \textbf{Trong khi chạy -- giao tiếp:} AGENT-RADAR quyết định agent chú ý vào câu nào trong lịch sử; DAIN học cường độ kết nối giữa các agent; UrbanAgent xây đồ thị quan hệ (corroboration, rectification, supervisory) giữa các agent.
  \item \textbf{Khi phơi bày kết quả:} MIRAGE tách ``execution view'' khỏi ``exposure view'' để đánh lừa kẻ tấn công.
  \item \textbf{Đánh giá sau cùng:} TIDE là một khung đánh giá, theo dõi dòng chảy bằng chứng đúng và sai qua toàn bộ quá trình.
\end{itemize}

\begin{figure}[!htbp]
\centering
\resizebox{\linewidth}{!}{%
\begin{tikzpicture}[
  font=\footnotesize,
  stage/.style={draw,thick,rounded corners,minimum width=2.75cm,minimum height=1.0cm,align=center,fill=black!5,font=\bfseries\footnotesize},
  tag/.style={draw,rounded corners=1pt,align=center,font=\scriptsize,minimum width=2.55cm,inner sep=2pt},
  arr/.style={-{Stealth[length=2.5mm]},very thick,black!55}
]
\node[stage] (p0) at (0,0)    {Thiết kế\\vai trò \& đồ thị};
\node[stage] (p1) at (3.15,0) {Điều phối\\\& kích hoạt};
\node[stage] (p2) at (6.3,0)  {Giao tiếp\\\& lập luận};
\node[stage] (p3) at (9.45,0) {Tổng hợp\\\& phơi bày};
\node[stage] (p4) at (12.6,0) {Đánh giá\\\& kiểm toán};
\draw[arr] (p0)--(p1); \draw[arr] (p1)--(p2); \draw[arr] (p2)--(p3); \draw[arr] (p3)--(p4);

\node[tag,fill=gOrch!15,below=0.4cm of p0] (h) {\#5 HalluProp\\(pre-hoc)};
\node[tag,fill=gOrch!15,below=0.4cm of p1] (pr) {\#3 ProgRouter};
\node[tag,fill=gComm!15,below=0.15cm of pr] (d1) {\#7 DAIN (cổng)};
\node[tag,fill=gComm!15,below=0.4cm of p2] (ra) {\#8 AGENT-RADAR};
\node[tag,fill=gComm!15,below=0.15cm of ra] (d2) {\#7 DAIN (đồ thị)};
\node[tag,fill=gApp!20,below=0.15cm of d2] (u) {\#6 UrbanAgent};
\node[tag,fill=gApp!20,below=0.15cm of u] (m) {\#2 MAC-RRG};
\node[tag,fill=gSafe!15,below=0.4cm of p3] (mi) {\#1 MIRAGE};
\node[tag,fill=gSafe!15,below=0.4cm of p4] (t) {\#4 TIDE};
\end{tikzpicture}}
\caption{Vị trí can thiệp của từng nghiên cứu trong vòng đời một multi-agent system.}
\label{fig:lifecycle}
\end{figure}

\section{Bản chất của ``agent'' trong các nghiên cứu}

Một quan sát quan trọng khi đọc kỹ tám bài là thuật ngữ ``agent'' được dùng với nghĩa rất khác nhau (Bảng~\ref{tab:agentnature}). Chỉ năm bài (\#1, \#3, \#4, \#5, \#8) dùng \emph{LLM agent thực sự} trao đổi bằng ngôn ngữ tự nhiên. UrbanAgent dùng LLM agent ở giai đoạn đầu nhưng giai đoạn ``debate'' là một bộ tổng hợp nơ-ron được huấn luyện. DAIN gọi các mô-đun nơ-ron con là ``agent'', về bản chất là một MoE có trao đổi thông điệp giữa các expert. MAC-RRG gọi các mô-đun retrieval tất định là ``agent''. Sự khác biệt này ảnh hưởng trực tiếp đến khả năng so sánh và được thảo luận thêm ở Chương~\ref{chap:comparison}.

\begin{table}[!htbp]
\centering
\small
\caption{Bản chất của agent và phương thức giao tiếp trong từng nghiên cứu.}
\label{tab:agentnature}
\begin{tabularx}{\textwidth}{L{2.4cm} L{3.3cm} Y L{2.4cm}}
\toprule
\textbf{Nghiên cứu} & \textbf{Agent là gì} & \textbf{Phương thức giao tiếp} & \textbf{Cần huấn luyện?} \\
\midrule
\#1 MIRAGE & LLM theo vai trò & Văn bản trên đồ thị DAG & Không \\
\#3 ProgRouter & LLM (vai trò thực thi) & Coordinator giao việc qua sổ cái (ledger) & Bộ dự báo học trực tuyến \\
\#4 TIDE & LLM (5 agent đồng nhất) & Phát quảng bá công khai + kênh riêng, 3 vòng & Không \\
\#5 HalluProp & LLM theo vai trò & Văn bản trên đồ thị, 3 vòng & Risk predictor có giám sát \\
\#8 AGENT-RADAR & LLM theo vai trò & Văn bản, đồ thị ngẫu nhiên/phân lớp/đầy đủ/debate & Không \\
\#6 UrbanAgent & VLM dùng công cụ (mỗi phương thức một agent) & Đồ thị quan hệ học được (truyền vector) & Có (SFT + GRPO + GNN) \\
\#7 DAIN & Mô-đun nơ-ron & Vector qua low-rank bottleneck & Có (đầu-cuối) \\
\#2 MAC-RRG & Mô-đun retrieval tất định & Ghép token tri thức vào tiền tố LLM & Có (LLM sinh báo cáo) \\
\bottomrule
\end{tabularx}
\end{table}
